\begin{pastexam}[label={px:ntu-109-dsaa-13}]{臺大資工所 109 學年度 第 13 題}
（問答）

(15 points) A disjoint-set data structure supports two operations. One is UNION($x,y$), which
merges the two roots of the trees containing $x$ and $y$. The other is FIND($x$), which returns
the root of the tree containing $x$. Union-by-height is a union heuristic, which keeps track of
the heights of the trees to attach the shorter tree to the root of the taller tree.
\begin{enumerate}[label=(\alph*)]
\item (9 points) Given an undirected graph $G=(V,E)$, where $V=\{v_1,\ldots,v_m\}$ ($m>0$) and
$E=\{e_1,\ldots,e_n\}$ ($n>0$), devise an algorithm to check if the graph $G$ is a tree or not by
using UNION($v_i,v_j$) with union-by-height and FIND($v_k$) with no path compression, where
$v_i,v_j$ and $v_k\in V$. Clearly describe your solution and analyze the time complexity of your
algorithm in the worst case. \textbf{No code is required (code will not be graded).}
\item (6 points) The union-by-height heuristic prevents the height of a tree from growing
linearly. Consider another heuristic, called union-by-descendant, which always attaches the tree
with fewer descendants to the root of the tree with more descendants. Suppose that path
compression is not applied. Which of the two heuristics is asymptotically better? Why?
\end{enumerate}
\end{pastexam}

\begin{pastsoln}
答案：(a) 先檢查 $n = m - 1$，再逐條邊用 FIND 看兩端是否已經在同一個集合（是就有迴路，不是樹），不是就 UNION；
最壞 $O(m \log m)$。 (b) 漸近一樣好，兩者都保證 $N$ 個元素時樹高 $O(\log N)$，而且都會達到 $\Theta(\log N)$。

這一題的記號和資結書相反：$m = |V|$ 是頂點數、$n = |E|$ 是邊數，下面照題目的記號。 不做路
徑壓縮時，union-by-height 就是資結書第~13.2.2~節的依秩合併：每個根記的秩恰好是樹的
高度（用邊數算），一樣高時任選一邊當根、高度加一。

\textbf{(a)} 由資結書定理~21.3，$G$ 是樹的充分必要條件是「$G$ 沒有迴路，而且 $n = m - 1$」。
所以先檢查邊數，再用互斥集合檢查有沒有迴路：

\begin{pseudo}[unbreakable]{判斷無向圖是不是樹}
\Function{Is-Tree}{$V, E$}
    \If{$n \ne m - 1$}
        \State \Return \FALSE
    \EndIf
    \For{每個頂點 $v \in V$}
        \State \Call{Make-Set}{$v$}
    \EndFor
    \For{每條邊 $(v_i, v_j) \in E$}
        \State $r_i \gets$ \Call{Find}{$v_i$}
        \State $r_j \gets$ \Call{Find}{$v_j$}
        \If{$r_i = r_j$}
            \State \Return \FALSE \Comment{這條邊和之前的邊形成迴路}
        \EndIf
        \State \Call{Union}{$v_i, v_j$}
    \EndFor
    \State \Return \TRUE
\EndFunction
\end{pseudo}

\textbf{正確性}。 用歸納法證明：處理完前 $t$ 條邊、還沒有回傳時，(i) 每個集合恰好是子圖
$G_t = (V, \{e_1, \ldots, e_t\})$ 的一個連通元件；(ii) $G_t$ 沒有迴路。 $t = 0$ 時每個頂點自
己一組、沒有邊，成立。 處理 $e_{t+1} = (v_i, v_j)$：
\begin{itemize}
\item 第 9 行成立：由 (i)，$v_i$、$v_j$ 在 $G_t$ 裡有一條路徑，加上 $e_{t+1}$ 就是 $G$ 裡的
一個迴路（$v_i = v_j$ 時 $e_{t+1}$ 本身是自環），$G$ 不是樹，回傳 \FALSE{} 正確。
\item 第 9 行不成立：$v_i$、$v_j$ 在 $G_t$ 的不同元件，加上 $e_{t+1}$ 正好把這兩個元件連成
一個，第 11 行也正好合併這兩個集合，(i) 仍然成立。 $e_{t+1}$ 不在任何迴路上（否則去掉它，
$G_t$ 裡還有一條 $v_i$ 到 $v_j$ 的路徑），(ii) 也成立。
\end{itemize}
迴圈跑完都沒回傳，表示 $G$ 沒有迴路；第 2 行又確認了 $n = m - 1$，由資結書定理~21.3
的 (3)$\Rightarrow$(1)，$G$ 是樹。 反過來，如果 $G$ 是樹，$n = m - 1$ 通過第 2 行，沒有任何邊
會讓第 9 行成立，最後回傳 \TRUE。

\textbf{時間}。 依秩合併不壓縮時，由資結書引理~13.3 (3)，每棵樹的高度不超過
$\lfloor \lg m \rfloor + 1$，所以每次 FIND 是 $O(\log m)$，UNION 裡面的兩次 FIND 也是。 第 2 行
$O(1)$；$m$ 次 \proc{Make-Set} 是 $O(m)$；第 6--11 行最多跑 $n = m - 1$ 次，每次 $O(\log m)$。
最壞情況的總時間是 $O(m \log m)$。 如果不先檢查邊數、改成最後數集合個數，迴圈要跑 $n$ 次，
總時間是 $O(m + n \log m)$；先檢查邊數可以在 $n$ 很大時立刻回答。

\textbf{(b)} 兩個規則在漸近意義下一樣好，都不比對方好。 根的子孫個數等於樹的節點數減一，所以「子孫
比較少的接到子孫比較多的底下」就是資結書第~13.2.1~節的依大小合併。 設一共有 $N$ 個元
素：
\begin{itemize}
\item 依大小合併：由資結書定理~13.2，有 $s$ 個節點的樹高度不超過
$\lfloor \lg s \rfloor + 1$。 理由是一個節點的深度加一時，它所在的樹被接到一棵至少一樣大的
樹底下，節點數至少加倍，所以深度最多 $\lg N$。
\item 依高度合併：由資結書引理~13.3 (2)，高度（邊數）是 $h$ 的樹至少有 $2^h$ 個節點：
兩棵一樣高的樹合併，高度變成 $h + 1$，節點數至少 $2^h + 2^h = 2^{h+1}$；不一樣高時高度不
變、節點數變多。 所以高度最多 $\lg N$。
\end{itemize}
兩個規則都保證 FIND 最壞 $O(\log N)$。 而且兩個都會達到 $\Theta(\log N)$：$N = 2^k$，每一輪
都把兩棵一模一樣的樹合併，$k$ 輪之後兩個規則做出同一棵樹，高度（邊數）$k = \lg N$，最深
的節點做 FIND 要走 $\lg N$ 步（資結書第~13.3~節 $N = 8$ 的例子）。 所以兩者最壞都是
FIND $\Theta(\log N)$、$M$ 個操作 $\Theta(M \log N)$，漸近相同。 對個別的操作序列，兩者做出
的樹高度可能不同（資結書第~13~章觀念題第~5~題就有依大小合併比較高的例子），但只差常數
倍，不影響漸近的結論。
\end{pastsoln}
